\section{About Score Functions: Sigmoid vs Softmax}
\label{app:scores}

\subsection{Vanilla CP Setting}

We recall the following score definitions: 

\begin{equation}
\begin{aligned}
    s_\mathrm{LAC \ Softmax} (x, y) &= 1 - \text{softmax} \left( \frac{f(x)}{T} \right)_y \\
    s_\mathrm{LAC \ Sigmoid} (x, y) &= 1 - \text{sigmoid} \left( \frac{f(x)_y - b}{T} \right)
\end{aligned}
\end{equation}

To evaluate the impact of using our LAC Sigmoid score as compared to a classic LAC Softmax score function, we provide the following empirical results for vanilla CP on the \textsl{CIFAR-10} dataset with a ResNet50 model with $\alpha = 0.1$:

\begin{table}[h]
  \centering
  \sisetup{table-format=2.2, detect-weight=true, detect-family=true}
  \begin{tabular}{l S S}
    \toprule
    Method       & {Coverage (\%)} & {Set size} \\
    \midrule
    LAC Softmax  & 90.04           & 1.088      \\
    LAC Sigmoid  & 90.28           & 1.144      \\
    \bottomrule
  \end{tabular}
  \caption{Comparison of conformal score performance LAC Softmax vs LAC Sigmoid.}
  \label{tab:scores}
\end{table}

This ablation study results in marginally bigger vanilla CP set sizes for LAC Sigmoid scores compared to softmax ones with scaled temperature.


\subsection{Robust CP Setting}

Using the 1-Lipschitz nature of our models, we can also use exact worst-case LAC Sigmoid and Softmax score computations under adversarial logit perturbations. Instead of relying on Eq.~\ref{eq:lipschitz_ineq}, we can leverage the monotonicity of the sigmoid function to compute tighter bounds, e.g:

\begin{equation}
    1 - \text{sigmoid}\left( \frac{f(\Tilde{x})_y + L_n \times \epsilon - b} {T} \right) \leq \min_{x \in \bouleps(\Tilde{x})} 
    s_\mathrm{LAC \ Sigmoid} (x, y)
    %\ \forall \Tilde{x} \in \bouleps(x).
\end{equation}

where $T$ and $b$ are the temperature and the bias of the score. This inequality yields tighter bounds than those based solely on global Lipschitz constants, with no extra cost.


Similarly, given that the softmax function exhibits monotonicity in each of its parameters (ceteris paribus), its maximum and minimum values within a hyper-rectangular region are attained at the corners of that region, i.e.:

\begin{equation}
    %\underline{s}(\Tilde{x}, y) = 
    1 - \text{softmax}
    \begin{pmatrix}
        (f(\Tilde{x})_0 - L_n  \times \epsilon) / T \\
        \vdots \\
        (f(\Tilde{x})_y + L_n \times \epsilon)/T \\
        \vdots \\
        (f(\Tilde{x})_{c-1} - L_n \times \epsilon) / T
    \end{pmatrix}_y \leq \min_{x \in \bouleps(\Tilde{x})} s_\mathrm{LAC \ Softmax}(x,y) % \forall \Tilde{x} \in \bouleps(x).
\end{equation}

where $T$ is the temperature scaling factor of the softmax.

As shown below, for a 1-Lipschitz VGG with 25M parameters on \textsl{CIFAR-10} with $\epsilon=0.03$ and $\alpha = 0.1$, LAC Sigmoid demonstrates considerably smaller robust prediction sets.

\begin{table}[h]
  \centering
  \sisetup{
    table-format=2.1,      % for Coverage
    table-number-alignment = center,
    detect-weight = true,
    detect-family = true
  }
  \begin{tabular}{l S S}
    \toprule
    Method        & {Coverage (\%)} & {Set size} \\
    \midrule
    LAC Softmax   & 95.6            & 2.38       \\
    LAC Sigmoid   & 94.9            & 1.72       \\
    \bottomrule
  \end{tabular}
  \caption{Robust CP performance for different conformal scores on \textsl{CIFAR-10}. Under the same conditions as Fig~\ref{fig:fastcp_results}.}
  \label{tab:robust_cp}
\end{table}

Although the LAC Sigmoid score is marginally worse in vanilla CP conditions, the superior performances obtained in Table~\ref{tab:robust_cp} motivates its use in the paper.
